% !TeX root = ../main.tex

\chapter{国内外研究现状}
\label{cha:related-work}

本章从推荐算法的整体演进出发，依次梳理大语言模型增强推荐、知识图谱增强推荐以及图采样与多智能体方法三个方向的研究进展，并分析现有工作的局限性，为本文第三章与第四章的研究动机提供依据。

\section{推荐算法研究概述}
\label{sec:rec-overview}

\subsection{从浅层模型到深度模型}
\label{subsec:shallow-to-deep}

早期推荐算法以协同过滤为基础，矩阵分解将用户与物品映射到低维隐空间，以向量内积刻画交互偏好\cite{koren2009mf}；因子分解机通过隐向量内积建模二阶特征交互\cite{rendle2010fm}。随着深度学习的发展，Wide\&Deep 将线性模型的记忆能力与深度网络的泛化能力相结合\cite{cheng2016wide}，DeepFM 进一步以因子分解机替代 Wide 部分实现端到端训练\cite{guo2017deepfm}，深度交叉网络则显式建模高阶特征交叉\cite{wang2017dcn}。

\subsection{用户兴趣建模}
\label{subsec:interest-modeling}

用户兴趣具有多样性与动态演化特性。DIN 引入注意力机制，根据候选物品自适应地为历史行为分配权重\cite{zhou2018din}；DIEN 在此基础上以门控循环单元建模兴趣随时间的演化过程\cite{zhou2019dien}。此类方法显著提升了兴趣刻画的细粒度，但注意力计算与序列建模带来了较高的在线推理开销。

\subsection{双塔架构与工业部署}
\label{subsec:two-tower}

为满足大规模在线服务的时延要求，双塔架构将用户与物品分别编码为独立向量，通过近似最近邻检索完成召回\cite{huang2013dssm}。SimpleX 等研究表明，经过恰当设计的双塔模型在协同过滤任务上同样具备强竞争力\cite{mao2021simplex}。双塔的瓶颈在于用户侧与物品侧缺乏交互，难以刻画细粒度兴趣。

\section{大语言模型增强的推荐方法}
\label{sec:llm4rec}

\subsection{基于提示与生成的推荐}
\label{subsec:llm-prompt}

该类方法将推荐任务转化为自然语言任务，通过提示工程直接让大语言模型生成推荐结果或解释\cite{zhao2024llmrec}。其优势在于无需训练即可利用开放世界知识，但存在推理成本高、结果不稳定、难以与既有推荐系统对接等问题。

\subsection{基于表征增强的推荐}
\label{subsec:llm-representation}

该类方法利用大语言模型生成用户与物品的语义表示，并将其注入传统推荐模型。LLM2DIN 以语言模型生成的物品嵌入初始化 DIN 模型\cite{harte2023llm2din}；KAR 利用大语言模型抽取用户偏好知识与物品事实知识，通过多专家编码器整合到 DeepFM 与 DIN 中\cite{xi2024kar}；LLM-ESR 面向长尾序列推荐，结合冻结的语义嵌入与协同信号进行双视角建模\cite{liu2024llmesr}；AlphaRec 则揭示了语言空间与行为空间之间的同态关系\cite{sheng2024alpharec}。本文第三章的工作属于这一路线。

\subsection{现有方法的局限}
\label{subsec:llm-limitation}

现有方法多将大语言模型的输出作为静态特征一次性注入，缺乏表示空间对齐的显式约束；同时，语义知识的引入往往伴随显著的推理开销，难以满足工业级时延要求。此外，多数方法面向序列推荐，对以点击率预测为目标的细粒度兴趣聚合关注不足。

\section{知识图谱增强的推荐方法}
\label{sec:kg4rec}

\subsection{基于知识图谱嵌入的方法}
\label{subsec:kg-embedding}

该类方法将知识图谱中的实体与关系嵌入低维向量空间，并作为物品侧特征引入推荐。DKN 将新闻文本与知识实体表示相结合实现知识感知的新闻推荐\cite{wang2018dkn}；MKR 通过多任务学习在推荐与知识图谱补全之间共享知识\cite{wang2019mkr}；RippleNet 模拟用户兴趣在知识图谱上的波纹式传播\cite{wang2018ripplenet}。

\subsection{基于图神经网络的方法}
\label{subsec:kg-gnn}

该类方法直接在用户—物品—实体构成的异构图上游走与聚合。KGCN 以感受野机制发现高阶结构信息\cite{wang2019kgcn}；KGAT 以注意力机制在协同知识图上传播信息\cite{wang2019kgat}；KGIN 进一步引入意图粒度建模关系路径\cite{wang2021kgin}。本文第四章的工作与该类方法关系最为紧密。

\subsection{现有方法的局限}
\label{subsec:kg-limitation}

现有方法普遍依赖人工构建或预先给定的知识图谱，图结构一经构建即静态不变，难以适应业务数据的动态演化；同时，长尾节点因交互稀疏而表征学习不充分，头部节点的冗余交互又引入噪声。

\section{图采样与多智能体方法}
\label{sec:sampling-agent}

图神经网络在推荐中的效果高度依赖于邻域采样策略。SGL 通过节点丢弃、边丢弃与随机游走构造自监督视图\cite{wu2021sgl}，验证了采样策略对表征质量的影响。另一方面，多智能体框架借助大语言模型的规划与工具调用能力，将复杂任务分解为若干子任务协同完成，已在知识构建与推理场景中展现出潜力。将多智能体引入图构建与采样过程，有望实现图结构的自动化、动态化维护。

\section{本章小结}
\label{sec:related-summary}

综上所述，大语言模型增强推荐与知识图谱增强推荐分别从语义与结构两个侧面注入外部知识，但前者面临表示对齐与推理效率的挑战，后者面临图结构静态与长尾稀疏的挑战。本文第三章与第四章分别针对上述两类问题展开研究，并在方法设计上有意形成“向量级知识注入”与“图结构级知识注入”的互补关系。
